\chapter{Pruebas}\label{cap:pruebas}

\section{Estrategia de pruebas}

La validación de este sistema debía combinar pruebas funcionales manuales, verificación del
flujo de datos entre frontend y backend y comprobaciones de usabilidad sobre la interfaz.
Dado que el proyecto depende de hardware real, la estrategia de pruebas debe contemplar
tanto escenarios controlados de laboratorio como pruebas de uso cotidiano.

En un sistema como GlucoMate no bastaba con verificar funciones aisladas. Es necesario
comprobar una cadena completa que incluye hardware, conectividad local, backend, base de
datos y renderizado en interfaz. Por ello, la estrategia más realista para este TFG es la
de prueba funcional de integración, complementada con comprobaciones manuales sobre cada
módulo: lectura de glucosa, registro de eventos, persistencia y asistencia por IA.

\section{Pruebas funcionales con sensor real Freestyle Libre}

Las pruebas funcionales más relevantes son las realizadas con un sensor real y un teléfono
Android con soporte NFC. La secuencia de validación adecuada para este proyecto puede
resumirse en los siguientes pasos:

\begin{enumerate}
    \item Abrir la aplicación móvil y verificar la carga del histórico, el perfil y los
          eventos clínicos recientes. \textbf{EXITOSA}
    \item Activar el flujo de escaneo NFC, aproximar el teléfono al sensor y que se detecte comprobando así el uso del chip NFC. \textbf{EXITOSA}
    \item Confirmar que el frontend obtiene una lectura válida desde el puente Juggluco. \textbf{EXITOSA}
    \item Verificar que la lectura se almacena correctamente en el backend y aparece en el
          histórico de 12 horas. \textbf{EXITOSA}
    \item Registrar una comida, una dosis de insulina y una actividad deportiva, y
          comprobar su reflejo inmediato en la interfaz. \textbf{EXITOSA}
    \item Enviar una consulta al agente de IA y validar que la respuesta utiliza el
          contexto clínico reciente. \textbf{EXITOSA}
\end{enumerate}

Este enfoque permitió comprobar no solo la funcionalidad de cada módulo, sino también la
coherencia temporal entre ellos, que es un aspecto esencial del sistema.

\section{Resultados y evidencias}

Los resultados obtenidos durante las pruebas confirman que el flujo principal de
GlucoMate se ejecuta correctamente de extremo a extremo. La aplicación es capaz de recibir
las lecturas procedentes del sensor, mostrarlas en la interfaz y mantenerlas sincronizadas
con el resto de módulos del sistema. En las comprobaciones realizadas, los valores de
glucosa recibidos por la aplicación coinciden con los valores que se pueden visualizar en
la aplicación oficial Abbott FreeStyle LibreLink, lo que permite validar que la lectura
integrada en GlucoMate conserva la información original del sensor sin alterar su valor
clínico ni su marca temporal.

Esta verificación se ha realizado comparando manualmente la lectura mostrada en GlucoMate
con la lectura visible en LibreLink para el mismo instante de medición. La coincidencia
entre ambas aplicaciones demuestra que el sistema interpreta correctamente los datos
recibidos y que el frontend representa de forma fiel la información capturada. Además, la
gráfica histórica y la lectura actual mantienen una evolución coherente, sin saltos
temporales injustificados ni discrepancias entre el dato puntual mostrado al usuario y el
histórico almacenado.

La persistencia en base de datos también se comporta de forma correcta. Cada lectura,
evento de comida, dosis de insulina, actividad física o registro clínico se almacena con
una marca temporal asociada, permitiendo reconstruir la secuencia de acontecimientos de
manera ordenada. Las pruebas realizadas evidencian que los datos guardados conservan
coherencia temporal entre sí y que, tras cerrar y volver a abrir la aplicación, la
información continúa disponible y se muestra de forma consistente en el frontend. Por
tanto, no se trata únicamente de una visualización temporal en la interfaz, sino de una
persistencia efectiva en el backend.

Las evidencias principales de este apartado se apoyan en tres elementos: capturas de la
pantalla principal con la lectura actual y la gráfica histórica, comprobaciones de la base
de datos con los registros persistidos y comparación directa con la aplicación oficial
Abbott FreeStyle LibreLink. En conjunto, estas evidencias permiten afirmar que el sistema
funciona de forma estable, que los datos se guardan correctamente y que existe coherencia
entre el sensor, la base de datos y la información presentada al usuario en la aplicación.

